 \subsection{\tt allocate\_block\_memory}
 This subroutine assigns to every block {\sl i} a number of node 
 points {\sl block(i)\%np} and a number of cells {\sl block(i)\%ncell}  
 The number of vertices per cell is also stored in {\sl block(i)\%nv}.
 It then loops over all blocks and allocates the required arrays 
 to store the mesh nodes position in Cartesian and spherical coordinates
 and the connectivity.  
 \subsection{\tt block\_node\_layout}
 This subroutine generates the position of the mesh node points for each 
 block as well as their connectivity (i.e. the list of nodes making 
 each element. 
 \begin{center}
 \includegraphics[width=8cm]{images/basics}\\
 {\small a) block layout for mtype=3,4; b) block layout for mtype=5}
 \end{center}
 In the case that the block is a quadrilateral, it is assumed that 
 it is made of $l \times l$ quadrilateral cells. In the case that the block is a 
 triangle, it is made of $l^2$ triangular cells.
 \subsection{\tt build\_hollow\_sphere}
 This subroutine first makes a copy of the current shell into {\sl shell\_temp}
 It then computes the number of mesh nodes and cells for the hollow sphere
 and allocate all arrays accordingly.
 It then loops over layers, places the temporary shell at the required 
 radius, and appends it to the current hollow sphere mesh from the inside out.
 \subsection{\tt check\_parameters}
 \subsection{\tt }
 \subsection{\tt hexahedron\_volume}
 This function computes the volume of any hexahedron, 
 following Eq.(12) of \cite{gran97}.
 \subsection{\tt compute\_gravity\_at\_point}
 This subroutine receives as argument the coordinates {\tt x,y,z} of a point and 
 returns the gravity vector components {\tt gx,gy,gz} at this location 
 by means of 
 \[
 {\bm g}({\bm r}') = \int_\Omega {\cal G} \frac{\rho({\bm r})-\rho_0}{|{\bm r}'-{\bm r}|^3} ({\bm r}'-{\bm r}) d\Omega
 \] 
 where $\rho_0=${\tt refdensgrav}.
 The gravitational potential is computed as follows:
 \[
 U({\bm r}') = -\int_\Omega {\cal G} \frac{\rho({\bm r})-\rho_0}{|{\bm r}'-{\bm r}|}  d\Omega
 \]
 \subsection{\tt compute\_gravity\_on\_line}
 This subroutine computes on a line parametrised by $0\leq r \leq 2R_2$,
 $\theta=13^\circ$ and $\phi=17^\circ$. This line is discretised 
 over npts points. At each point is the subroutine compute\_gravity\_at\_point called.
 \subsection{\tt compute\_hollow\_sphere\_volume}
 \subsection{\tt compute\_normals\_hollow\_sphere}
 \subsection{\tt compute\_r\_theta\_phi\_hollow\_sphere}
 This subroutine computes $r$, $\theta$ and $\phi$ from $x,y,z$ for all 
 nodes of the hollow\_sphere mesh with
 \[
 r=\sqrt{x^2+y^2+z^2}
 \]
 \[
 \theta=\cos^{-1}(z/r)
 \]
 \[
 \phi=\tan^{-1}_2 (y/x)
 \]
 \subsection{\tt compute\_r\_theta\_phi\_shell}
 This subroutine computes $r$, $\theta$ and $\phi$ from $x,y,z$ for all 
 nodes of the shell mesh with
 \[
 r=\sqrt{x^2+y^2+z^2}
 \]
 \[
 \theta=\cos^{-1}(z/r)
 \]
 \[
 \phi=\tan^{-1}_2 (y/x)
 \]
 \subsection{\tt compute\_shell\_area}
 \subsection{\tt compute\_shell\_distributions}
 \subsection{\tt compute\_triangle\_area}
 {\tt https://en.wikipedia.org/wiki/Heron\%27s\_formula }
 Heron's formula states that the area of a triangle whose sides have lengths a, b, and c is
 \[
 A=\sqrt{s(s-a)(s-b)(s-c)}
 \]
 where s is the semiperimeter of the triangle; that is,
 \[
 s=\frac{a+b+c}{2}
 \]
 \subsection{\tt generate\_quadrature\_points}
 \begin{center}
 \includegraphics[width=6cm]{images/elements}
 \end{center}
 For hexahedral elements the linear shape functions are:
 \begin{eqnarray}
 N_1(r,s,t) &=& (1-r)(1-s)(1-t)/8 \nonumber\\ 
 N_2(r,s,t) &=& (1+r)(1-s)(1-t)/8 \nonumber\\ 
 N_3(r,s,t) &=& (1+r)(1+s)(1-t)/8 \nonumber\\ 
 N_4(r,s,t) &=& (1-r)(1+s)(1-t)/8 \nonumber\\ 
 N_5(r,s,t) &=& (1-r)(1-s)(1+t)/8 \nonumber\\ 
 N_6(r,s,t) &=& (1+r)(1-s)(1+t)/8 \nonumber\\ 
 N_7(r,s,t) &=& (1+r)(1+s)(1+t)/8 \nonumber\\ 
 N_8(r,s,t) &=& (1-r)(1+s)(1+t)/8 \nonumber
 \end{eqnarray}
 For the triangular prism elements the linear shape functions are:
 \begin{eqnarray}
 N_1(r,s,t) &=& (1-r-s)(1-t)/2 \nonumber\\ 
 N_2(r,s,t) &=& r(1-t)/2 \nonumber\\ 
 N_3(r,s,t) &=& s(1-t)/2 \nonumber\\ 
 N_4(r,s,t) &=& (1-r-s)(1+t)/2 \nonumber\\ 
 N_5(r,s,t) &=& r(1+t)/2 \nonumber\\ 
 N_6(r,s,t) &=& s(1+t)/2 \nonumber
 \end{eqnarray}
 \subsection{\tt map\_blocks}
 \subsubsection{HS06 - cubed sphere}
 See section 3.1 of \cite{puli07}
 \subsubsection{HS12 - CitcomS mesh}
 \subsubsection{HS20 - Icosahedron}
 Let $t=(1+\sqrt{5})/2$, then the vertices of the icosahedron are 
 \[
 {\bm x}_0=\frac{1}{\sqrt{1+t^2}}(t,1,0) \quad\quad {\bm x}_6=\frac{1}{\sqrt{1+t^2}}(-1,0,t)
 \]
 \[
 {\bm x}_1=\frac{1}{\sqrt{1+t^2}}(-t,1,0) \quad\quad {\bm x}_7=\frac{1}{\sqrt{1+t^2}}(-1,0,-t)
 \]
 \[
 {\bm x}_2=\frac{1}{\sqrt{1+t^2}}(t,-1,0) \quad\quad {\bm x}_8=\frac{1}{\sqrt{1+t^2}}(0,t,1)
 \]
 \[
 {\bm x}_3=\frac{1}{\sqrt{1+t^2}}(-t,-1,0) \quad\quad {\bm x}_9=\frac{1}{\sqrt{1+t^2}}(0,-t,1)
 \]
 \[
 {\bm x}_4=\frac{1}{\sqrt{1+t^2}}(1,0,t) \quad\quad {\bm x}_{10}=\frac{1}{\sqrt{1+t^2}}(0,t,-1)
 \]
 \[
 {\bm x}_5=\frac{1}{\sqrt{1+t^2}}(1,0,-t) \quad\quad {\bm x}_{11}=\frac{1}{\sqrt{1+t^2}}(0,-t,-1)
 \]
 It is obvious that all points lie on a sphere of unit radius.

 The triangles are as follows:
 \[  T_0=[0,8,4]  \quad\quad T_{10}=[2,9,11]   \]
 \[  T_1=[0,5,10] \quad\quad T_{11}=[3,11,9]   \]
 \[  T_2=[2,4,9]  \quad\quad T_{12}=[4,2,0]   \]
 \[  T_3=[2,11,5] \quad\quad T_{13}=[5,0,2]   \]
 \[  T_4=[1,6,8]  \quad\quad T_{14}=[6,1,3]   \]
 \[  T_5=[1,10,7] \quad\quad T_{15}=[7,3,1]   \]
 \[  T_6=[3,9,6]  \quad\quad T_{16}=[8,6,4]   \]
 \[  T_7=[3,7,11] \quad\quad T_{17}=[9,4,6]   \]
 \[  T_8=[0,10,8] \quad\quad T_{18}=[10,5,7]   \]
 \[  T_9=[1,8,10] \quad\quad T_{19}=[11,7,5]   \]
 Graphically, the mesh has the following structure:
 \begin{center}
 \includegraphics[width=12cm]{images/flat_icosahedron}
 \end{center}
 \subsection{\tt merge\_blocks}
 \subsection{\tt output\_blocks\_vtu}
 \subsection{\tt output\_hollow\_sphere\_vtu}
 \subsection{\tt output\_shell\_ascii}
 \subsection{\tt output\_shell\_lat\_lon}
 This subroutine exports $\theta$ and $\phi$ values of all points on the 
 shell in the file {\sl shell\_theta\_phi.dat} in {\sl OUT}.
 \subsection{\tt output\_shell\_vtu}
 If {\sl generate\_vtu\_output} is true 
 a vtu file of the shell is created in {\sl OUT}.
 \subsection{\tt project\_on\_sphere}
 This subroutine gets a radius $r$ value as argument and the current coordinates
 of a point. It returns the new coordinates of the point on a sphere 
 of radius $r$ with the same $\theta,\phi$.
 \subsection{\tt read\_command\_line\_options}
 This subroutine reads in the command line arguments
 \subsection{\tt read\_s40rts}
 \subsection{\tt }
 \subsection{\tt write\_positions}
 \subsection{\tt write\_field\_dp}
 \subsection{\tt write\_field\_int}
 \subsection{\tt write\_icon}
 \subsection{\tt write\_offsets}
 \subsection{\tt write\_types}
